% theory/ParabolicDiffusion/GEAR_FVPM_Hyperbolic_Diffusion_Theory.tex

\documentclass{article}
\usepackage[margin=1in]{geometry} % Wider text block: this document is dense
                                   % with long inline \code{} identifiers and
                                   % display equations that overflow the
                                   % default article textwidth.
\usepackage{amsmath}
\usepackage{amssymb}
\usepackage{amsthm}
\usepackage{graphicx}
\usepackage{hyperref}
\usepackage[utf8]{inputenc} % Required for Unicode characters
\usepackage{xcolor}         % Required for coloring text
\usepackage{enumitem}       % For customizing list environments
\usepackage{physics}        % For common physics notations like dv, grad, div, etc.
\usepackage{booktabs}       % For clean tables
\usepackage{listings}       % For inline/blocked source-code identifiers

% ---------------------------------------------------------------------------
% Code-identifier typesetting. \lstinline treats underscores literally, so
% there is no risk of an unescaped "_" breaking the build.
% ---------------------------------------------------------------------------
\lstset{
  basicstyle=\ttfamily\small,
  breaklines=true,
  columns=fullflexible,
  keepspaces=true
}
\newcommand{\code}[1]{\mbox{\lstinline{#1}}}

\theoremstyle{plain}
\newtheorem{proposition}{Proposition}
\newtheorem{lemma}{Lemma}
\newtheorem{corollary}{Corollary}
\theoremstyle{definition}
\newtheorem{definition}{Definition}
\theoremstyle{remark}
\newtheorem{remark}{Remark}

% ---------------------------------------------------------------------------
% Physics macros, shared with GEAR_FVPM_Diffusion_Theory.tex (parabolic
% companion document) for notational continuity, plus a few new ones needed
% only for the hyperbolic (relaxation) system.
% ---------------------------------------------------------------------------
\newcommand{\Zfrac}{Z}                                  % metal mass fraction
\newcommand{\rhoZ}{\rho \Zfrac}                          % metal mass density
\newcommand{\driver}{q}                                  % diffusion driver
\newcommand{\Kmat}{\mathbf{K}}                            % diffusion tensor
\newcommand{\Kstar}{\mathbf{K}^{\star}}                   % interface diffusion tensor
\newcommand{\Shear}{\mathbf{S}}                           % shear tensor
\newcommand{\Splus}{\mathbf{S}_{+}}                       % regularised (PSD) shear tensor
\newcommand{\Fdiff}{\mathbf{F}}                           % diffusive flux (algebraic in \S1, dynamical from \S2 on)
\newcommand{\kappaC}{\kappa_{0}}                          % chem\_data->diffusion\_coefficient
\newcommand{\kappaP}{\kappa}                              % per-particle chemistry\_data.kappa
\newcommand{\CCFL}{C_{\mathrm{CFL}}}                      % chem\_data->C\_CFL\_chemistry
\newcommand{\gammak}{\gamma_{k}}                          % kernel\_gamma
\newcommand{\nunit}{\mathbf{n}}                           % interface unit normal
\newcommand{\dx}{\Delta x}                                % GIZMO particle size
\newcommand{\dt}{\Delta t}                                % explicit timestep
\newcommand{\Uvec}{\mathbf{U}}                            % hyperbolic state vector
\newcommand{\Gflux}[1]{\mathbf{G}^{(#1)}}                 % flux tensor, direction-#1 component
\newcommand{\chyp}{c_{\rm hyp}}                           % hyperbolic diffusion sound speed
\newcommand{\taurelax}{\tau}                              % relaxation time (unmacroed in the parent doc; kept explicit here)

\title{Wave-Speed Theory for the Hyperbolic (Cattaneo-Relaxation) FVPM
Diffusion Scheme}
\author{Darwin Roduit}
\date{\today}

\begin{document}
\sloppy % several long, unhyphenatable \code{} identifiers would otherwise
        % overfull the line; prefer looser justification to clipped text.
\maketitle
\tableofcontents

\section{Scope and Motivation}
\label{sec:scope}

This is the companion document promised, but not written, at the end of
Section~\emph{Hyperbolic Relation} of
\emph{Theoretical Foundations of Parabolic Diffusion in GEAR FVPM}
(\code{GEAR_FVPM_Diffusion_Theory.tex}, this directory): a first-principles
derivation of the characteristic (wave) speed of the hyperbolic
(Cattaneo~\cite{Cattaneo1948}-relaxation) diffusion system implemented under
\code{CHEMISTRY_GEAR_FVPM_HYPERBOLIC_DIFFUSION}
(\code{src/chemistry/GEAR_FVPM_DIFFUSION/hyperbolic/}), and a direct
confrontation of that result against the two closed-form expressions
currently returned by
\code{chemistry_get_physical_hyperbolic_soundspeed}
(\code{hyperbolic/chemistry_getters.h:37--66}).

The motivating question was operational: \emph{does the hyperbolic
timestep still need a $\dt<\taurelax$-type restriction, given that the
relaxation source term is integrated exactly?} Answering it precisely
required first pinning down what the \emph{correct}, general (mode- and
shape-of-$\Kmat$-independent) characteristic speed actually is --
independent of any of the three closed forms already in the code, and
without presupposing which of them (if any) is right. That derivation is
the content of this document. Several conclusions followed directly
enough from Proposition~\ref{prop:main} to implement in the same pass:
coupling $\taurelax_0$ to $\kappaC$ (Section~\ref{sec:tau0-physical}), and
fixing the \code{constant_mode} soundspeed branch's missing $\rho^{-1}$
and its use of the Frobenius norm in place of a tight, still
particle-local bound (Remark~\ref{rem:lambda-max-not-n}). What remains
open is stated as findings and recommendations in
Section~\ref{sec:confrontation}, left for a deliberate decision. This is
part of the ongoing hyperbolic-diffusion verification work on the
\code{darwin/gear_chemistry_fluxes} branch.

\subsection{Relation to the Parabolic Document}
This document assumes familiarity with the FVPM geometry, the diffusion
driver $\driver$, and the diffusion tensor $\Kmat$ as defined in
Sections~2--3 of the parabolic companion document; it does not repeat that
material. The two documents share notation (see the macro list at the top
of the source) so that symbols mean the same thing in both.

\section{Governing System}
\label{sec:governing-system}

\subsection{Conservation of Metal Mass}
As in the parabolic case (parabolic document, Eq.~1), conservation of
metal-mass density in the frame comoving with each Lagrangian particle
reads
\begin{equation}
  \label{eq:mass_conservation}
  \pdv{(\rhoZ)}{t} + \div \Fdiff = 0 .
\end{equation}
The difference between the parabolic and hyperbolic schemes is entirely in
what closes this equation: the parabolic scheme slaves $\Fdiff$
algebraically to $\grad\driver$ via Fick's law (parabolic document,
Eq.~3); the hyperbolic scheme instead promotes $\Fdiff$ to an independent
dynamical variable, evolved by its own PDE.

\subsection{The Cattaneo Relaxation Equation}
The flux obeys (\code{chemistry_flux.h:90--94}, docstring of
\code{chemistry_part_integrate_flux_source_term}):
\begin{equation}
  \label{eq:cattaneo}
  \pdv{\Fdiff}{t} + \frac{\Kmat}{\taurelax}\grad\driver = -\frac{\Fdiff}{\taurelax} ,
\end{equation}
where $\taurelax = \code{p->chemistry_data.tau}$ is the local relaxation
time (Section~\ref{sec:tau-modes}) and $\driver$ is the same diffusion
driver as in the parabolic scheme ($\driver=\rhoZ$ for
\code{isotropic_constant}; $\driver=\Zfrac$ otherwise -- see
\code{chemistry_compute_flux}, \code{hyperbolic/chemistry_flux.h:159--172}).
As $\taurelax\to0$, Eq.~\eqref{eq:cattaneo} forces
$\Fdiff\to-\Kmat\grad\driver$, i.e.\ it collapses onto Fick's law and
Eq.~\eqref{eq:mass_conservation} becomes the parabolic equation -- this is
the relaxation limit already noted (without proof) in the parabolic
document.

The pair \eqref{eq:mass_conservation}--\eqref{eq:cattaneo} is exactly the
Cattaneo/Maxwell-Cattaneo regularisation of a diffusion law, structurally
identical to the M1 closure of radiative transfer
(Berthon et al.~\cite{Berthon2007}) -- the same structural identity already
noted informally in the code's development log and reflected in the
blending scheme of Section~\ref{sec:confrontation}.

\subsection{Relaxation-Time Models}
\label{sec:tau-modes}
Two independent modes exist
(\code{chemistry_compute_physical_tau}, \code{hyperbolic/chemistry_getters.h:75--92}):
\begin{align}
  \text{\code{constant_mode}:} \qquad & \taurelax = \taurelax_0 , \\
  \text{\code{turbulent_mode}:} \qquad & \taurelax = \frac{\taurelax_0}{\kappaC\,\norm{\Shear}_F} ,
  \label{eq:tau_turbulent}
\end{align}
with $\taurelax_0=$\code{chem_data->tau} the parameter-file normalisation
and $\Shear$ the (unregularised) physical velocity shear tensor. The
\code{turbulent_mode} formula includes an explicit $\kappaC^{-1}$
(\code{chem_data->diffusion_coefficient}) coupling, motivated and derived
in Section~\ref{sec:tau0-physical}; it supersedes an earlier,
$\kappaC$-independent version of this formula (see that section for why
the coupling is required). This choice of relaxation-time model is
orthogonal, in the parameter file, to the diffusion-tensor mode $\Kmat$
(Section~\ref{sec:diffusion-tensor-recap}): \code{diffusion_mode} and
\code{relaxation_time_mode} are parsed as two independent strings
(\code{chemistry.h:279--326}), with no cross-validation between them.

\subsection{Diffusion Tensor Modes (recap)}
\label{sec:diffusion-tensor-recap}
From \code{chemistry_get_physical_matrix_K}/\code{chemistry_get_diffusion_coefficient}
(\code{chemistry_getters.h:196--285}):
\begin{align}
  \text{\code{isotropic_constant}:} \qquad & \Kmat=\kappaP\mathbf I, \quad \kappaP=\kappaC, \\
  \text{\code{isotropic_smagorinsky}:} \qquad & \Kmat=\kappaP\mathbf I, \quad \kappaP=\kappaC\,\gammak^2\,h^2\,\rho\,\norm{\Shear}_F, \\
  \text{\code{anisotropic_gradient}:} \qquad & \Kmat = \kappaC\,\gammak^2\,h^2\,\rho \; \Splus .
\end{align}
Crucially, $\Kmat$ carries an explicit factor of $\rho$ in the
\code{isotropic_smagorinsky} and \code{anisotropic_gradient} modes, but
\emph{not} in \code{isotropic_constant}. This asymmetry is the origin of
the apparent inconsistency examined in Section~\ref{sec:confrontation}.

\section{Quasi-Linear Form and the Flux Jacobian}
\label{sec:quasilinear}

\subsection{State Vector and Frozen Coefficients}
Define the four-component hyperbolic state vector
\begin{equation}
  \Uvec = (U_0, U_1, U_2, U_3) \equiv (\rhoZ, F_x, F_y, F_z) .
\end{equation}
As is standard for deriving a local wave-speed estimate for a Godunov-type
solver (and as the code itself already does implicitly, by evaluating
$\chyp$ from each particle's own, locally stored $\Kmat,\taurelax,\rho$ --
\code{chemistry_iact.h:295--298}), we treat $\Kmat$, $\taurelax$ and $\rho$
as \emph{locally frozen} coefficients for the purpose of constructing the
flux Jacobian: they are evaluated at the particle/interface, not
differentiated with respect to $\Uvec$. This is the same frozen-coefficient
approximation used, without comment, by every scheme in this code family
(GIZMO's own MFV/MFM sound-speed estimate is exactly analogous). Under this
approximation $\driver = U_0/\rho$ is linear in $U_0$ (both for
$\driver=\rhoZ=U_0$ trivially, and for $\driver=\Zfrac=U_0/\rho$ with
$\rho$ frozen).

\subsection{The Flux Tensor}
Directly from the code (\code{chemistry_get_hyperbolic_flux},
\code{hyperbolic/chemistry_getters.h:162--200}), the flux of $\Uvec$ in
Cartesian direction $k\in\{x,y,z\}$ is
\begin{equation}
  \Gflux{k}(\Uvec) = \left( F_k, \; \frac{K_{xk}\driver}{\taurelax}, \; \frac{K_{yk}\driver}{\taurelax}, \; \frac{K_{zk}\driver}{\taurelax} \right) ,
\end{equation}
so that Eqs.~\eqref{eq:mass_conservation}--\eqref{eq:cattaneo} read, in
conservative form, $\partial_t\Uvec + \sum_k \partial_{x_k}\Gflux{k}(\Uvec)
= -(0,\Fdiff/\taurelax)$. This is an exact transcription of the
implementation, not an independent modelling choice: row 0 of
\code{hypflux} is literally \code{U[1..3]}, and rows 1--3 are literally
\code{K[j][k]*(q/tau)} (\code{hyperbolic/chemistry_getters.h:178--199}).

\subsection{The Directional Flux Jacobian}
For a Riemann problem at an interface with unit normal $\nunit$, the
relevant $4\times4$ matrix is the projected flux Jacobian
$\mathbf A(\nunit) \equiv \partial(\Gflux{}(\Uvec)\cdot\nunit)/\partial\Uvec$.
Writing $\mathbf v \equiv \Kmat\nunit/(\rho\taurelax)$ (a 3-vector), direct
differentiation of the flux tensor above gives, for any $\mathbf x =
(x_0,x_1,x_2,x_3)$,
\begin{equation}
  \label{eq:jacobian_action}
  \big(\mathbf A(\nunit)\,\mathbf x\big)_0 = \nunit\cdot(x_1,x_2,x_3) ,
  \qquad
  \big(\mathbf A(\nunit)\,\mathbf x\big)_j = v_j\, x_0 \quad (j=1,2,3) .
\end{equation}

\section{The Characteristic Speed: General Result}
\label{sec:main-theorem}

\begin{proposition}[Characteristic speeds of the hyperbolic diffusion system]
\label{prop:main}
For any symmetric positive semi-definite $3\times3$ diffusion tensor
$\Kmat$ (no assumption on its shape -- isotropic, Smagorinsky-scaled, or
fully anisotropic), any relaxation-time model, and any unit interface
normal $\nunit$, the projected flux Jacobian $\mathbf A(\nunit)$
of Eq.~\eqref{eq:jacobian_action} has eigenvalues
\begin{equation}
  \label{eq:main_result}
  \lambda \in \left\{ \; +\chyp(\nunit),\; -\chyp(\nunit),\; 0,\; 0 \; \right\} ,
  \qquad
  \boxed{\;\chyp(\nunit) = \sqrt{\dfrac{\nunit^{\!\top}\Kmat\,\nunit}{\rho\,\taurelax}}\;} .
\end{equation}
\end{proposition}

\begin{proof}
Let $\lambda\neq0$ be an eigenvalue with eigenvector $\mathbf x$. From the
second block of Eq.~\eqref{eq:jacobian_action}, $\lambda x_j = v_j x_0$ for
$j=1,2,3$, i.e.\ $x_j = (v_j/\lambda)\,x_0$ (this requires $x_0\neq0$: if
$x_0=0$ then $x_j=0$ for all $j$ from this relation, i.e.\ $\mathbf
x=\mathbf 0$, not a valid eigenvector). Substituting into the first block,
$\lambda x_0 = \nunit\cdot(x_1,x_2,x_3) = (\nunit\cdot\mathbf v/\lambda)\,x_0$,
and since $x_0\neq0$,
\begin{equation}
  \lambda^2 = \nunit\cdot\mathbf v = \frac{\nunit\cdot(\Kmat\nunit)}{\rho\taurelax} = \frac{\nunit^{\!\top}\Kmat\nunit}{\rho\taurelax} \; \geq 0
\end{equation}
(non-negative because $\Kmat$ is PSD), giving the two roots
$\lambda=\pm\chyp(\nunit)$ of Eq.~\eqref{eq:main_result}. The remaining two
eigenvalues are $\lambda=0$: any $\mathbf x=(0,x_1,x_2,x_3)$ with
$\nunit\cdot(x_1,x_2,x_3)=0$ (a two-dimensional subspace, the plane
transverse to $\nunit$) satisfies $\mathbf A(\nunit)\mathbf x=\mathbf 0$
by direct substitution into Eq.~\eqref{eq:jacobian_action}.
\end{proof}

\begin{remark}
The two zero eigenvalues correspond to the flux components transverse to
the interface normal: they are advected with zero characteristic speed in
this frozen-coefficient, one-dimensional-projection sense (they do not
"propagate" through the Riemann fan). This is the same structure as the
transverse magnetic field components in a 1D MHD Riemann problem, or the
transverse momentum components in a 1D gas-dynamics Riemann problem: real,
but not part of the acoustic-family wave-speed bound. The code's HLL
solver (Section~\ref{sec:confrontation}) uses the same $\lambda_\pm$ bound
for all four flux components (\code{chemistry_riemann_HLL.h:112-179}),
which is the standard, conservative simplification -- not a further
approximation beyond what Proposition~\ref{prop:main} already licenses for
the two genuinely propagating families.
\end{remark}

Proposition~\ref{prop:main} was derived using only
Eqs.~\eqref{eq:mass_conservation}--\eqref{eq:cattaneo} and the
frozen-coefficient flux Jacobian -- no assumption was made anywhere about
the internal structure of $\Kmat$ (isotropic, Smagorinsky, or
anisotropic). The two corollaries below specialise it.

\subsection{Corollary: Isotropic Recovery}
\begin{corollary}
\label{cor:isotropic}
If $\Kmat=\kappaP\mathbf I$ (both \code{isotropic_constant} and
\code{isotropic_smagorinsky} modes), then $\nunit^{\!\top}\Kmat\nunit =
\kappaP$ for \emph{every} unit vector $\nunit$, and
\begin{equation}
  \chyp = \sqrt{\frac{\kappaP}{\rho\,\taurelax}} ,
\end{equation}
independent of direction and \emph{independent of whether $\kappaP$ came
from the constant or the Smagorinsky closure.}
\end{corollary}
This is the formal statement of what the user's own by-hand derivation
anticipated: the factor of $\rho^{-1}$ is required by the general
Proposition, not by any assumption on the shape of $\Kmat$ -- it appears
identically whether $\Kmat$ is isotropic-constant or isotropic-Smagorinsky.
It is a consequence of $\driver$ being related to the \emph{conserved}
state $U_0=\rhoZ$ by $\driver=U_0/\rho$, not of how $\kappaP$ itself was
computed.

\subsection{Corollary: The Frobenius Norm Is a (Non-Tight) Upper Bound}
\begin{corollary}
\label{cor:frobenius}
For any unit $\nunit$ and any symmetric PSD $\Kmat$ with eigenvalues
$\lambda_1\geq\lambda_2\geq\lambda_3\geq0$,
\begin{equation}
  \nunit^{\!\top}\Kmat\nunit \; \leq \; \lambda_1 \; \leq \; \norm{\Kmat}_F ,
  \qquad
  \frac{\norm{\Kmat}_F}{\lambda_1} \in [1,\sqrt3\,] ,
\end{equation}
with the upper end of the ratio, $\sqrt3$, attained \emph{only} in the
isotropic case $\lambda_1=\lambda_2=\lambda_3$, and the lower end, $1$,
attained only for a rank-1 (maximally anisotropic) $\Kmat$.
\end{corollary}
\begin{proof}
Write $\nunit=\sum_i c_i\mathbf e_i$ in the eigenbasis of $\Kmat$
($\sum_ic_i^2=1$); then $\nunit^{\!\top}\Kmat\nunit=\sum_i\lambda_ic_i^2
\leq\lambda_1\sum_ic_i^2=\lambda_1$, with equality iff $\nunit=\mathbf
e_1$. Also $\norm{\Kmat}_F^2=\sum_i\lambda_i^2\geq\lambda_1^2$, with
equality iff $\lambda_2=\lambda_3=0$ (rank 1), so $\norm{\Kmat}_F/\lambda_1
\geq1$. For the upper end, in 3D $\sum_i\lambda_i^2 \leq
3\max_i\lambda_i^2=3\lambda_1^2$ with equality iff all three eigenvalues
are equal, giving $\norm{\Kmat}_F/\lambda_1\leq\sqrt3$.
\end{proof}
Concretely, evaluating $\norm{\Kmat}_F$ for $\Kmat=\kappaP\mathbf I$ gives
$\norm{\Kmat}_F=\kappaP\sqrt3$, \emph{not} $\kappaP$: using the Frobenius
norm in place of the direction-projected $\nunit^{\!\top}\Kmat\nunit$
overestimates the true characteristic speed by a factor of exactly
$\sqrt3\approx1.73$ even in the isotropic limit, and by a smaller (but
never negative) factor for anisotropic $\Kmat$. It is always a
\emph{safe} (CFL-stable) overestimate, never an underestimate, but it is
also never tight except for rank-1 $\Kmat$.

\begin{remark}[$\lambda_1(\Kmat)$, not $\nunit$, is the right practical fix]
\label{rem:lambda-max-not-n}
Proposition~\ref{prop:main} is stated per-interface, with $\chyp(\nunit)$
depending on the pair-specific normal $\nunit$. It is tempting to conclude
that fixing $\chyp$ in the code requires threading $\nunit$ into
\code{chemistry_get_physical_hyperbolic_soundspeed} and evaluating it
per-interaction rather than once per particle -- but this is unnecessary,
and would be a worse engineering choice. What the CFL condition and the
HLL wave-speed bound actually need is not the exact directional value, but
\emph{any} bound that is never smaller than $\nunit^{\!\top}\Kmat\nunit$
for every $\nunit$ a particle might face. Corollary~\ref{cor:frobenius}
already proves $\lambda_1(\Kmat)$ -- the largest eigenvalue of a
particle's own $\Kmat$ -- is exactly such a bound, and a strictly tighter
one than $\norm{\Kmat}_F$ (equal to it only for isotropic $\Kmat$; strictly
smaller, and exact for the worst-case direction, otherwise). $\lambda_1(\Kmat)$
depends only on $\Kmat$, hence only on particle-local quantities ($\rho$,
$h$, $\Shear$ or $\Splus$) -- it \emph{is} a particle property, computable
once per particle from an eigendecomposition of $\Kmat$ (SWIFT already has
one, \code{sym\_matrix\_diagonalise\_3x3\_d}, used elsewhere in this module
by \code{chemistry\_regularize\_shear\_tensor}), exactly like $\taurelax$,
$\kappaP$, or the old $\norm{\Kmat}_F$ it replaces. No call site needs a
new argument. The only thing given up relative to the fully general,
per-interface $\nunit^{\!\top}\Kmat\nunit$ is tightness for anisotropic
$\Kmat$ when the interface does \emph{not} happen to align with $\Kmat$'s
dominant eigendirection -- in that case $\lambda_1(\Kmat)$ is still a safe
overestimate, just not an exact one. Given Corollary~\ref{cor:frobenius}
already bounds how loose that overestimate can be
($\leq\sqrt3\,\lambda_1$, only reached when $\Kmat$ is isotropic -- in
which case $\lambda_1$ \emph{is} exact for every $\nunit$ regardless), this
is a good trade: it fixes both open gaps (the missing $\rho^{-1}$ and the
Frobenius looseness) with a local, per-particle change, not an
architectural one.
\end{remark}

\section{Physically-Motivated Normalisation of $\taurelax_0$}
\label{sec:tau0-physical}

The propositions above fix the \emph{shape} dependence of $\chyp$ but say
nothing about how large $\taurelax_0$ itself should be. Numerically,
$\taurelax_0$ is a free parameter-file constant; physically, in the
\code{turbulent_mode} closure it is supposed to represent \emph{how many
local shear/turnover times the diffusive flux takes to relax to its
Fickian equilibrium}. This section derives why that constant cannot be
chosen independently of $\kappaC$, checks the resulting formula against an
independent literature derivation, and revises the ``does hyperbolic beat
parabolic'' criterion of the earlier discussion accordingly. The
conclusions here are already implemented in
\code{hyperbolic/chemistry_getters.h} (both
\code{chemistry_compute_physical_tau} and
\code{chemistry_get_physical_hyperbolic_soundspeed}).

\subsection{Why $\taurelax_0$ Must Depend on $\kappaC$}
\label{sec:tau0-motivation}
Whether the hyperbolic scheme is faster than the parabolic one is decided
by comparing $\taurelax$ against the intrinsic parabolic crossing time of
the same problem, $\taurelax_{\rm par}\sim\dx^2/D$ -- exactly the ratio the
code already forms internally for the HLL blending factor
(\code{chemistry_riemann_compute_hyperbolic_blending_factor},
\code{chemistry_riemann_utils.h}). Since $D$ (through $\kappaP$) is itself
set by $\kappaC$, this comparison timescale is not independent of
$\kappaC$: doubling $\kappaC$ halves $\taurelax_{\rm par}$, so a
$\taurelax_0$ that does not track $\kappaC$ silently drifts to a different
point on the parabolic$\leftrightarrow$hyperbolic spectrum every time
$\kappaC$ is retuned (e.g.\ during the flux-limiter calibration of
Section~\ref{sec:confrontation}'s companion work). A physically meaningful
$\taurelax_0$ must therefore be defined \emph{relative to} $\kappaC$, not
alongside it.

\subsection{Cross-Check: The Parabolic Mixing Timescale of Romano et al.~(2022)}
Romano, Nagamine \& Hirashita~\cite{Romano2022} use the same
Shen-et-al.~\cite{ShenWadsleyStinson2010}/Smagorinsky-Lilly~\cite{Smagorinsky1963}
closure implemented here (their diffusivity, eq.~10 of that paper, is
$D=C_d\norm{\Shear_{ab}}h^2\rho$ -- identical in structure to our
$\kappaP$ for \code{isotropic_smagorinsky}, up to the SPH kernel-support
factor $\gammak^2$). They derive the physical timescale for turbulent
diffusion to erase a contrast $\Delta A$ in some transported quantity $A$
(their eq.~13):
\begin{equation}
  \label{eq:romano_tdiff}
  \taurelax_{{\rm diff},A} \sim \left(\frac{A}{\Delta A}\right)
  \left(\frac{C_d}{0.01}\right)^{-1}
  \left(\frac{\norm{\Shear_{ab}}}{100\ {\rm km\,s^{-1}\,kpc^{-1}}}\right)^{-1}\ {\rm Gyr} ,
\end{equation}
i.e.\ $\taurelax_{{\rm diff},A}\propto(\kappaC\norm{\Shear}_F)^{-1}$. This
is precisely the $\kappaC^{-1}\norm{\Shear}_F^{-1}$ scaling
Section~\ref{sec:tau0-motivation} argues is required -- an independent,
literature-derived confirmation that the coupling is physical, not merely
a numerical convenience. It also has no explicit dependence on resolution
$h$, for the same reason found in this document: $D\propto h^2$ cancels
the $h^2$ implicit in a diffusive crossing time $\dx^2/D$.

\subsection{Implemented Reparametrisation}
Equation~\eqref{eq:tau_turbulent} already states the adopted form,
$\taurelax=\taurelax_0/(\kappaC\norm{\Shear}_F)$. Propagating this through
$\chyp=\sqrt{\kappaP/(\rho\taurelax)}$ with the isotropic Smagorinsky
$\kappaP=\kappaC\gammak^2h^2\rho\norm{\Shear}_F$
(Corollary~\ref{cor:isotropic}) gives the revised closed form:
\begin{equation}
  \label{eq:code_turbulent_new}
  \chyp^{\rm code} = \frac{\kappaC}{\sqrt{\taurelax_0}}\,\gammak\,h\,\norm{\Shear}_F ,
\end{equation}
which replaces the pre-reparametrisation form
\begin{equation}
  \label{eq:code_turbulent_old}
  \chyp^{\rm code,\,old} = \sqrt{\frac{\kappaC}{\taurelax_0}}\,\gammak\,h\,\norm{\Shear}_F
\end{equation}
that this document originally derived (and that the code implemented)
before the $\kappaC$ coupling of Section~\ref{sec:tau0-motivation} was
adopted. Note the change from $\sqrt{\kappaC/\taurelax_0}$
(Eq.~\eqref{eq:code_turbulent_old}) to $\kappaC/\sqrt{\taurelax_0}$:
coupling $\taurelax$ to $\kappaC$ adds one further power of $\kappaC$ under
the square root in $\chyp^2=\kappaP/(\rho\taurelax)$, since $\kappaP\propto
\kappaC$ and now $1/\taurelax\propto\kappaC$ as well.
$\taurelax_0$ itself remains dimensionless (a pure count of shear times),
since $\kappaC$ is dimensionless in the Smagorinsky closure -- the
reparametrisation changes what value of $\taurelax_0$ corresponds to a
given physical relaxation time, but not its units. \textbf{This is a
breaking change for any existing calibrated $\taurelax_0$ in
\code{turbulent_mode}} (no example currently uses this mode, so nothing in
the repository needed updating).

\subsection{On Adding a $\driver/\Delta\driver$ Term}
\label{sec:no-q-over-dq}
A natural follow-up question: since Eq.~\eqref{eq:romano_tdiff} carries a
factor $A/\Delta A$, should $\taurelax$ carry an analogous
$\driver/\Delta\driver$ factor? \textbf{No.} $\taurelax_{{\rm diff},A}$
above is the timescale for the \emph{state} $A$ itself to relax
appreciably; it is a solution-dependent, emergent quantity that combines
the flux dynamics with how far the current state is from equilibrium.
$\taurelax$ in the Cattaneo system is a different object: a material/
turbulence property governing how fast the \emph{flux} responds to an
instantaneous local gradient, independent of how large the resulting
state change happens to be -- exactly analogous to a mean-free time in the
kinetic-theory derivation of Fourier's law, which does not depend on the
temperature contrast either. The $\driver/\Delta\driver$-like behaviour is
not missing from this formulation: it emerges automatically as a
\emph{consequence} of solving Eqs.~\eqref{eq:mass_conservation}--\eqref{eq:cattaneo}
together (in the $\taurelax\to0$ limit the system converges to the
parabolic equation, whose own relaxation time is exactly
Eq.~\eqref{eq:romano_tdiff}). Building $\driver/\Delta\driver$ directly
into $\taurelax$ would double-count this and would additionally break the
frozen-coefficient assumption of Section~\ref{sec:quasilinear} ($\taurelax$
would depend on $U_0$ and its gradient, invalidating the flux-Jacobian
derivation of Proposition~\ref{prop:main} as stated -- a real structural
cost, not just a modelling objection). Where a $\driver/\Delta\driver$-like
quantity legitimately belongs is in \emph{calibrating} the constant
$\taurelax_0$ itself: Eq.~\eqref{eq:romano_tdiff} can be evaluated at a
fiducial, order-unity contrast to translate a target physical mixing
timescale (in Gyr) into a concrete $\taurelax_0$ -- a calibration aid for
the constant, not a runtime dependence for $\taurelax$.

\begin{remark}[Ad hoc timestep criteria vs.\ a proper CFL condition]
Romano et al.~\cite{Romano2022} note that a na\"ive Courant-like criterion
on the parabolic update itself, $\dt\leq\alpha\,A/\lvert{\rm d}A/{\rm
d}t\rvert$ (their eq.~14), diverges as $A\to0$ (pristine gas), and work
around it with an ad hoc combination of an abundance floor
($A_{\rm low}=10^{-4}A_\odot$, below which outflows are unconstrained) and
asymmetric coefficients ($\alpha=5$ for inflows, $\alpha=0.5$ for
outflows), validated post-hoc by checking that a stricter choice
($\alpha=0.1$) does not change the results. This is precisely the class of
problem the hyperbolic reformulation is architecturally built to avoid:
because $\dt$ here is set by a genuine wave-speed CFL condition
($\dt=\CCFL\dx/v_{\max}$) rather than by the local relative rate of change
of the transported quantity, it does not need a bespoke floor, asymmetric
coefficients, or a post-hoc convergence check against a stricter variant
to handle the near-zero-abundance limit. The GEAR FVPM flux limiter's own
\code{startup_fraction} parameter addresses the same pristine-gas
pathology, but at the flux-limiting stage rather than the timestep stage,
and is ad hoc in the same sense (an empirically-set constant, not a
first-principles result) -- but it is deliberately \emph{not} pinned to a
fixed abundance scale the way $A_{\rm low}$ is: it caps the incoming flux
at a fraction (currently $10\%$) of the \emph{source} particle's own mass,
a dimensionless, scale-relative criterion that requires no problem- or
unit-system-specific floor to be chosen. It shares Romano et al.'s
reliance on an untuned constant, but not their reliance on an absolute
threshold for what counts as ``pristine.''
\end{remark}

\subsection{Revised Speed Criterion}
\label{sec:revised-speed-criterion}
Repeating the timestep-ratio argument with
Eq.~\eqref{eq:code_turbulent_new} in place of the code's previous
(uncoupled) formula, and using $v_{\max}\approx2\chyp$ (comparable
neighbouring particles, negligible relative velocity along $\nunit$ --
the same simplifying assumption used throughout this comparison):
\begin{equation}
  \dt_{\rm hyp} \approx \frac{\CCFL\dx}{2\chyp}
  = \frac{\CCFL\sqrt{\taurelax_0}}{2\,\kappaC\norm{\Shear}_F} ,
  \qquad
  \dt_{\rm par} = \frac{\CCFL}{\sqrt3\,\kappaC\norm{\Shear}_F}
\end{equation}
(the second expression is the code's actual \code{isotropic_smagorinsky}
parabolic timestep, \code{chemistry_timesteps.h:107--113}, which -- unlike
the now-corrected hyperbolic branch -- still uses the Frobenius norm
$\norm{\Kmat}_F=\sqrt3\,\kappaP$ rather than a direction-projected
quantity; that asymmetry is noted here but is out of scope to fix, since
Corollary~\ref{cor:frobenius} is a safe overestimate for a parabolic CFL
too). The resolution $h$, diffusivity normalisation $\kappaC$, and local
shear $\norm{\Shear}_F$ all cancel from the ratio:
\begin{equation}
  \label{eq:revised_ratio}
  \boxed{\;\frac{\dt_{\rm hyp}}{\dt_{\rm par}} = \frac{\sqrt3}{2}\sqrt{\taurelax_0} \;}
  \qquad\Longrightarrow\qquad
  \text{hyperbolic faster iff } \taurelax_0 \gtrsim \frac{4}{3} .
\end{equation}
This is a substantially cleaner criterion than the pre-reparametrisation
result (Section~\ref{sec:tau0-motivation}'s companion memory note; a
resolution- and shear-independent threshold requiring only knowledge of
$\kappaC$): after coupling $\taurelax_0$ to $\kappaC$, whether the
hyperbolic scheme is faster than parabolic no longer depends on
resolution, the diffusion normalisation, or the local shear at all -- only
on the single dimensionless constant $\taurelax_0$ itself, against a fixed
threshold of order unity.

\section{Confrontation with the Implementation}
\label{sec:confrontation}

\code{chemistry_get_physical_hyperbolic_soundspeed}
(\code{hyperbolic/chemistry_getters.h:37--66}) returns two different
closed forms depending on \code{relaxation_time_mode}, reproduced here for
comparison against Eq.~\eqref{eq:main_result}. The \code{turbulent_mode}
branch is Eq.~\eqref{eq:code_turbulent_new} of
Section~\ref{sec:tau0-physical} (already the current, $\kappaC$-coupled
form); the \code{constant_mode} branch, following
Remark~\ref{rem:lambda-max-not-n}, is now:
\begin{align}
  \text{\code{constant_mode} branch (implemented):} \qquad
  \chyp^{\rm code} &= \sqrt{\frac{\lambda_1(\Kmat)}{\rho\,\taurelax}} , \label{eq:code_constant}
\end{align}
which replaces the pre-fix form
$\chyp^{\rm code,\,old}=\sqrt{\norm{\Kmat}_F/\taurelax}$ that this
document originally found missing $\rho^{-1}$ and using the
non-tight Frobenius norm.

\begin{table}[h]
\centering
\footnotesize
\begin{tabular}{@{}p{4.2cm}p{6.5cm}p{4cm}@{}}
\toprule
\code{diffusion\_mode} $\times$ \code{relaxation\_time\_mode} & Code formula vs.\ Prop.~\ref{prop:main} & Verdict \\
\midrule
\code{isotropic\_constant} + \code{constant\_mode} & Eq.~\eqref{eq:code_constant}: $\lambda_1=\kappaP$ for isotropic $\Kmat$, $\equiv\sqrt{\kappaP/(\rho\taurelax)}$ & \textbf{consistent (exact)} \\
\code{isotropic\_smagorinsky} + \code{constant\_mode} & same as above, $\kappaP$ now carries $\rho$ & \textbf{consistent (exact)} \\
\code{anisotropic\_gradient} + \code{constant\_mode} & Eq.~\eqref{eq:code_constant} uses $\lambda_1(\Kmat)$, a tight bound on $\nunit^{\!\top}\Kmat\nunit$ (Cor.~\ref{cor:frobenius}) & \textbf{consistent (safe bound, exact for the worst-case $\nunit$)} \\
\code{isotropic\_smagorinsky} + \code{turbulent\_mode} & Eq.~\eqref{eq:code_turbulent_new} $\equiv$ Cor.~\ref{cor:isotropic} exactly & \textbf{consistent} \\
\code{isotropic\_constant} + \code{turbulent\_mode} & Eq.~\eqref{eq:code_turbulent_new} assumes Smagorinsky $\kappaP(h,\rho,\Shear)$ & formula does not apply; silently wrong \\
\code{anisotropic\_gradient} + \code{turbulent\_mode} & Eq.~\eqref{eq:code_turbulent_new} assumes isotropic $\Kmat$ & formula does not apply; silently wrong \\
\bottomrule
\end{tabular}
\caption{Every reachable (\code{diffusion\_mode}, \code{relaxation\_time\_mode})
combination (the two settings are parsed independently, \code{chemistry.h:279--326},
with no cross-validation) against the general result of Proposition~\ref{prop:main}.
Four of six pairings are now consistent, up from one; the two
\code{turbulent\_mode} mismatches (rows 5--6) remain open, since fixing
them requires the \code{turbulent_mode} branch to use $\Kmat$ and
$\lambda_1(\Kmat)$ directly instead of its hand-derived,
Smagorinsky-specific closed form.}
\end{table}

\subsection{Reading the Table}
\begin{itemize}
\item \textbf{The \code{constant_mode} branch is now consistent for every
  diffusion mode} (Remark~\ref{rem:lambda-max-not-n}): using
  $\lambda_1(\Kmat)$ with the $\rho^{-1}$ factor fixes the two gaps this
  document originally found there in one, purely per-particle change --
  no interface normal needed. This was the more consequential of the two
  original gaps ($\rho$ varies by orders of magnitude across a
  cosmological simulation; the superseded $\sqrt3$ Frobenius gap was only
  ever a bounded, $O(1)$ constant), and is now resolved. The
  currently-tested example
  (\code{examples/ChemistryTests/MetalDiffusionOnePeak/params.yml}:
  \code{isotropic_constant} + \code{Constant} tau) sits in the first row
  of the table and is exact after this fix (isotropic $\Kmat$: $\lambda_1=
  \kappaP$ for every interface direction).
\item \textbf{The \code{turbulent_mode} formula remains mode-specific, not
  general}: it still silently returns a value derived under the
  assumption $\Kmat=\kappaP\mathbf I$ with the particular Smagorinsky
  $\kappaP$, even when \code{diffusion_mode} is set to something else.
  Nothing prevents that combination being selected in the parameter file.
  Unlike the \code{constant_mode} fix, resolving this needs the
  \code{turbulent_mode} branch to stop hand-deriving a Smagorinsky-only
  closed form and instead call
  \code{chemistry_get_physical_matrix_K}/\code{chemistry_get_matrix_max_eigenvalue}
  directly, the same way the now-fixed \code{constant_mode} branch does --
  a small, still purely-particle-local change, but not yet made.
\end{itemize}

\section{Consequence for the HLL Numerical Dissipation}
\label{sec:hll-consequence}

An overestimated $\chyp$ does not only affect the timestep -- it directly
inflates the numerical dissipation of the (unlimited) hyperbolic HLL flux
used when the blending factor $\alpha\to1$
(\code{chemistry_riemann_HLL.h:112--133}; see the companion memory note on
the large-$\taurelax$ over-diffusivity gap). With
$\lambda_+=-\lambda_-\equiv\Lambda$ (the code's symmetric-ish estimate,
\code{chemistry_riemann_HLL.h:74--75}), the HLL dissipation term simplifies
to
\begin{equation}
  F_{\rm diss} = \frac{\lambda_+\lambda_-\,(U_R-U_L)}{\lambda_+-\lambda_-}
  = \frac{-\Lambda^2(U_R-U_L)}{2\Lambda} = -\frac{\Lambda}{2}(U_R-U_L) ,
\end{equation}
i.e.\ \emph{linear} in the wave-speed estimate $\Lambda$. Any factor by
which $\chyp$ is overestimated (the missing $\rho^{-1}$ where relevant, the
$\sqrt3$ Frobenius gap, or both) passes through directly into extra
numerical diffusion in the large-$\taurelax$ ($\alpha\approx1$)
regime -- on top of, and independent from, the missing Hopkins-2017-type
limiter on $F_{\rm diss}$ already identified in that regime. Tightening
$\chyp$ (using $\nunit^{\!\top}\Kmat\nunit$ rather than $\norm{\Kmat}_F$,
and adding the missing $\rho^{-1}$) is therefore a second, independent
lever on the ``too diffusive at large $\taurelax$'' symptom, complementary
to porting the Hopkins-2017 limiter itself.

\section{Summary of Findings}
\label{sec:summary}

Three changes have been applied to source as a direct consequence of this
document (items 2, 3 and 6 below); the remaining finding (4) is offered
for a deliberate decision, not applied automatically.

\begin{enumerate}
\item \textbf{Proposition~\ref{prop:main} is the general result}: $\chyp(\nunit)
  = \sqrt{\nunit^{\!\top}\Kmat\nunit/(\rho\taurelax)}$, derived from the
  conservation and Cattaneo-relaxation equations alone, with no assumption
  on the shape of $\Kmat$. This confirms the premise of the request that
  motivated this document: a mode-independent derivation is possible, and
  it is the $\rho^{-1}$ factor, not the choice of relaxation-time model,
  that is the universal, non-negotiable piece.
\item \textbf{Implemented.} The \code{constant_mode} branch of
  \code{chemistry_get_physical_hyperbolic_soundspeed} was missing the
  $\rho^{-1}$ factor, for every \code{diffusion_mode} -- the higher-priority
  of the two original gaps (unbounded in magnitude, unlike the next item).
  Now fixed: it divides by $\rho$ (\code{hydro\_get\_physical\_density}).
\item \textbf{Implemented, via Remark~\ref{rem:lambda-max-not-n}, not by
  threading $\nunit$.} The \code{constant_mode} branch used $\norm{\Kmat}_F$
  in place of the direction-projected $\nunit^{\!\top}\Kmat\nunit$, an
  overestimate of $\chyp$ by a bounded factor $\in[1,\sqrt3]$
  (Corollary~\ref{cor:frobenius}). Rather than making $\chyp$ a
  per-interaction quantity (which would need $\nunit$ threaded through
  \code{chemistry_iact.h} and \code{chemistry_riemann_HLL.h}, and would
  turn a once-per-particle computation into a per-neighbour one), it now
  uses $\lambda_1(\Kmat)$, the largest eigenvalue of the particle's own
  $\Kmat$ (new helper \code{chemistry\_get\_matrix\_max\_eigenvalue}, built
  on the existing \code{sym\_matrix\_diagonalise\_3x3\_d}) -- a strictly
  tighter, still purely particle-local bound: exact for isotropic $\Kmat$
  (every reachable \code{diffusion\_mode} paired with \code{constant\_mode}
  is now exact, not merely bounded), and a tighter-than-Frobenius safe
  overestimate for anisotropic $\Kmat$.
\item \emph{(Not yet implemented.)} The \code{turbulent_mode} branch's
  closed form (Eq.~\eqref{eq:code_turbulent_new}) is still only valid for
  \code{isotropic_smagorinsky} $\Kmat$; nothing in the code prevents
  pairing it with \code{anisotropic_gradient} or \code{isotropic_constant},
  where it would return a value not licensed by Proposition~\ref{prop:main}.
  Fixing this does not need $\nunit$ either -- the same
  $\lambda_1(\Kmat)/(\rho\taurelax)$ approach as item 3 applies -- it is a
  smaller, still-open piece of work, not a different kind of fix.
\item Findings 2--3 (now fixed) used to inflate the HLL
  numerical-dissipation term linearly (Section~\ref{sec:hll-consequence});
  finding 4 (still open) still does, for the \code{anisotropic_gradient}
  and \code{isotropic_constant} pairings with \code{turbulent_mode}. Not
  purely a timestep-accounting issue -- plausibly still contributes to the
  large-$\taurelax$ over-diffusivity symptom for those pairings,
  independently of the missing Hopkins-2017 limiter.
\item \textbf{Implemented.} $\taurelax_0$ in \code{turbulent_mode} is now
  coupled to $\kappaC$ (Section~\ref{sec:tau0-physical}:
  $\taurelax=\taurelax_0/(\kappaC\norm{\Shear}_F)$), matching the physical
  mixing-timescale scaling independently derived by
  Romano et al.~\cite{Romano2022}
  (Eq.~\eqref{eq:romano_tdiff}). Both \code{chemistry_compute_physical_tau}
  and \code{chemistry_get_physical_hyperbolic_soundspeed} were updated
  consistently. No example parameter file used \code{turbulent_mode} prior
  to this change, so the break in the meaning of an existing calibrated
  $\taurelax_0$ (Section~\ref{sec:tau0-physical}) affects no committed
  example. A $\driver/\Delta\driver$-type term was considered and
  deliberately \emph{not} added to $\taurelax$ itself
  (Section~\ref{sec:no-q-over-dq}): it would double-count physics that
  already emerges from solving the coupled system, and it would break the
  frozen-coefficient assumption underpinning Proposition~\ref{prop:main}.
\item All three source changes (items 2, 3, 6) were verified by rebuilding
  \code{src/chemistry.o} against the tree's configured
  \code{GEAR-FVPM-HYPERBOLIC-DIFFUSION\_10} variant with
  \code{-Wall~-Wextra~-Werror}: clean, no warnings.
\end{enumerate}

\bibliographystyle{plain}
\bibliography{references}

\end{document}
